% =====================================================================================
\chapter{EXPERIMENTAL SETUP AND PROCEDURE}
\label{chp:setup}
% Second shape 2026-10-20 (review B1b sections A-C: the pose server, the printed tool and the
% parts, the calibration tools). Shape of the examples (Oydu ch. 4: robot, software, hardware,
% test environment, procedure). Values are those in use on 2026-10-20.
% Facts the repository does not hold (the holder's printer and material, the pen, the payload
% setting) are marked ASK; data that lives on the desktop are marked DESKTOP.
% =====================================================================================

This chapter describes the hardware and software of the cell, the test parts, how the mark error
is measured, and how one experiment is run; the last section describes the setup of the WeldSet
benchmark. Chapter~\ref{chp:results} reports the results.

% -------------------------------------------------------------------------------------
\section{Robot, Camera and Tool}
\label{sec:setup_robot}
Table~\ref{tab:hardware} lists the hardware. The UR5e carries a six-axis force/torque sensor at
its flange, read at 500~Hz; the cell uses only the force, for touch detection, collision aborts
and the tool-centre and table calibrations. The camera is an Intel RealSense D435i, mounted on the
wrist beside the pen; it streams colour and depth aligned to colour at $1280\times720$ and 30~Hz.
Its stereo baseline is 50~mm and its minimum depth at this resolution about 0.28~m, which sets the
closest view of the multi-view refinement.

\begin{table}[htbp]
\caption{Hardware of the cell.}
\label{tab:hardware}
\centering\small
\begin{tabular}{@{}L{0.17\textwidth}L{0.47\textwidth}L{0.28\textwidth}@{}}
\toprule
Component & Model & Role \\
\midrule
Manipulator & Universal Robots UR5e, factory-calibrated kinematics, flange F/T sensor & motion; force sensing \\
Camera & Intel RealSense D435i, eye-in-hand & RGB-D, $1280\times720$, 30~Hz \\
Tool & pen in a 3D-printed holder with the camera mount & marks in place of a torch \\
Desktop & NVIDIA RTX 5070 Ti & registration (SAM2, PPF, FoundationPose) \\
Laptop & NVIDIA RTX 4060 Laptop GPU, 8~GB & ROS~2, tracking, refinement, planning, robot interface \\
Table & the robot's table, measured as a tilted plane (\S\ref{sec:cal_table}) & the height reference \\
\bottomrule
\end{tabular}
\end{table}
\TODO{ASK: the CPU and RAM of both machines; the UR5e payload and centre-of-gravity setting used
with the tool; the controller software version.}

\paragraph{The pen holder and camera mount.} The tool was designed for this work and 3D printed
(Figure~\ref{fig:tool}). It is one assembly bolted to the flange: a flange plate with a side arm
that carries the camera, on the side opposite the robot's tool connector, and a holder body in
which two clamp pieces grip the pen. The assembly measures $94\times126\times140$~mm; the holder
body reaches from 60 to 140~mm above the flange face. The pen tip, with the cap off, sits at
$(1.07,\,-1.35,\,180.91)$~mm in the flange frame (the tool-centre calibration of
\S\ref{sec:cal_tcp}). The camera's optical axis is parallel to the pen within about $0.2^\circ$,
with its optical centre at about $(34.8,\,89.0,\,54.5)$~mm in the flange frame (the calibrated
extrinsic, Chapter~\ref{chp:results}). The collision envelope of the planner
(\S\ref{sec:mc_collision}) was derived from the assembly's CAD model with a 3~mm radial margin, and
the camera body was measured on the bench. The camera's cable runs round the wrist, which is why
the last wrist joint is kept within $\pm180^\circ$ of its home value (\S\ref{sec:mc_reach}).
\TODO{ASK: the printer and filament of the holder (the repository documents the Creality K1 Max
only for the test parts), the infill, the pen (make, tip diameter), and whether the clamp holds the
pen rigidly (the pen is unsprung: \S\ref{sec:cal_table}). Photograph of the tool.}
\figplaceholder{fig:tool}{The pen tool: (a) photograph on the robot; (b) the CAD assembly (flange
plate with camera arm, holder body, clamps); (c) its collision envelope in RViz with the calibrated
optical origin}{The pen tool and its collision envelope.}
\SRC{world/penholder\_assembly.step (FreeCAD, 2026-07-04) and .usda; config/pen\_tool.json;
scripts/tool\_model\_marker\_node.py; review B1b \S B.1.}

% -------------------------------------------------------------------------------------
\section{Software}
\label{sec:setup_software}
The laptop runs Ubuntu 24.04 with ROS~2 Jazzy and Python 3.12, the Universal Robots ROS~2 driver
(its scaled joint trajectory controller for every planned motion and its velocity controller for the
camera aim) and the cell's package, a mix of Python nodes and three C++ components
(\S\ref{sec:impl_cpp}). The registration and tracking servers run in Docker, in one image built for
each GPU architecture, with Ubuntu 22.04 and Python 3.10; the image holds FoundationPose, SAM~2.1 and
OpenCV's point-pair features. Table~\ref{tab:software} gives the versions.

\begin{table}[htbp]
\caption{Software and versions (laptop, 2026-10-20).}
\label{tab:software}
\centering\small
\begin{tabular}{@{}L{0.36\textwidth}L{0.24\textwidth}L{0.32\textwidth}@{}}
\toprule
Component & Version & Used for \\
\midrule
Ubuntu, ROS~2, Python & 24.04, Jazzy, 3.12.3 & the cell \\
Universal Robots ROS~2 driver & 3.8.0 & controllers, F/T, TCP pose \\
UR description (URDF) & 3.5.1 & robot model, with the factory calibration \\
MoveIt core (time-optimal parameterisation) & 2.12.4 & trajectory timing \\
OMPL & 1.7.0 & APS, RRT-Connect \\
pybind11 & 2.11.1 & the C++ planner extension \\
\texttt{depth\_image\_proc} & 5.0.12 & organized point cloud \\
MediaPipe & 0.10.35 & gesture recognition \\
NumPy, SciPy, trimesh, OpenCV & 2.3.5, 1.18.0, 4.12.2, 5.0.0 & geometry, ICP, calibration tools \\
Docker image (desktop / laptop) & DESKTOP & FoundationPose, SAM~2.1 Hiera-S, OpenCV \texttt{ppf\_match\_3d}, Flask, WebSocket \\
\bottomrule
\end{tabular}
\end{table}
\TODO{DESKTOP: the image's versions (CUDA, PyTorch, FoundationPose commit, SAM2 commit, OpenCV
contrib), its Dockerfile name (docker/Dockerfile.blackwell; the laptop build with
TORCH\_CUDA\_ARCH\_LIST=8.9) and the run scripts; they live in the FoundationPose repository on the
GPU hosts, not in this workspace.}

\paragraph{Network.} The desktop is reached over a private overlay network (Tailscale), and only one
HTTP request per registration crosses it; frames for tracking never leave the laptop
(\S\ref{sec:impl_arch}). On the university's network, which blocks the overlay, a VPN is needed
before registration.

% -------------------------------------------------------------------------------------
\section{Test Parts}
\label{sec:setup_parts}
The test parts cover the seam families of the cell (Table~\ref{tab:setup_parts},
Figure~\ref{fig:setup_parts}). The plates are cut from 8~mm MDF stock; the tubes, the
rounded-rectangle section and the band were 3D printed on a Creality K1 Max, each from one watertight
CAD model in millimetres, from which both the printed STL and the cell's PLY are exported. The
filament is matte and light: black, glossy or translucent filament gives the D435i holes or wrong
depth. Walls are at least 3~mm, tubes are printed upright and bands and plates lying flat, and the print
list prescribes measuring every part with callipers after printing. Every part is in the CAD registry, fitted to
its exact primitive (\S\ref{sec:mc_registry}).

\begin{table}[htbp]
\caption{Test parts and assemblies (nominal dimensions).}
\label{tab:setup_parts}
\centering\small
\begin{tabular}{@{}L{0.16\textwidth}L{0.47\textwidth}L{0.29\textwidth}@{}}
\toprule
Part & Geometry & Seam on its base \\
\midrule
Base plate & MDF, $220\times220\times5.5$~mm (before 2026-10-15: $256\times249\times8$~mm) & --- \\
Ear & MDF, $250\times99\times8$~mm, standing on the base & two straight fillets (T-joint) \\
C1 & tube $\varnothing62$~mm, wall 3, length 100, flat ends & a circle \\
E2 & tube $\varnothing90$~mm, wall 5, length 110, base mitred at $25^\circ$ & an ellipse \\
R2 & tube $\varnothing100$~mm, wall 6, length 180, lying on two V-blocks & (the main tube of S3) \\
S3 & tube $\varnothing50$~mm, wall 4, length 100, base cut to sit on R2 at $70^\circ$ & a saddle on R2 \\
RR1 & rounded-rectangle section $64\times64$~mm, corner radius 12, wall 3, height 100 & a closed rounded rectangle \\
SP3 & S-shaped band, wall 5, height 80, span about 220~mm (cubic B-spline spine) & two curved fillets \\
\bottomrule
\end{tabular}
\end{table}
\TODO{ASK: the filament material (PLA?), and the measured dimensions (callipers) next to the
nominal ones; the real MDF thickness (the README measured 8.4~mm for the ``8~mm'' ear); which base
plate each experiment used. State one description of E2 (a vertical tube with a $25^\circ$ cut, so
it stands tilted).}

\paragraph{Holding the parts.} The tubes and the band stand by themselves on the base plate; R2 lies
on two MDF V-blocks, and some parts are held by magnets, which can let them move by about a
millimetre between the scan, the marking and the touches. Fixtures are entered into the collision
model as a few pen-tip points on their tops (\S\ref{sec:mc_collision}).
\TODO{OPEN (todo.md): enter the V-blocks and holders as fixtures (one of SP3's first tacks hit an
unmodelled fixture).}
\figplaceholder{fig:setup_parts}{Photographs of the test assemblies on the table: the T-joint, C1,
E2, S3 on R2, RR1, SP3}{The test assemblies.}
\SRC{notes/curved\_parts\_print\_list.md (2026-10-08/09); models/*.ply,
models/weldgen\_objects.json; test/data/bench\_models/README.md; review B1b \S B.2.}

% -------------------------------------------------------------------------------------
\section{Measuring the Mark Error}
\label{sec:setup_measure}
The reference for every accuracy claim must not depend on the camera. Two measurements satisfy
this. The first is the pen's own contact (Equation~\ref{eq:touch}): every descent stops at the first
1.5~N touch, so the tip's position at that moment, from the robot's calibrated kinematics, is a point
on the real part, and its signed distance from the registered point along the pen axis, the contact
depth $\delta_c$, is the registration's error in that direction. Because a T-joint's two fillets are
touched along axes $45^\circ$ apart, early contact on both sides indicates a vertical error and early
contact on one side with late contact on the other a sideways shift; 1~mm of height error reads as
about 1.4~mm of contact depth along a $45^\circ$ pen. The second is the mark itself, a pen dot on the
part, whose distance from the true weld root is measured on the part.
\TODO{Define the second measurement: how the true root is found (the edge of the standing plate, a
scribed line), the instrument (ruler, callipers, a flatbed scan of the plate), its resolution, and
which claim uses which reference.}

Two further tools are used for evaluation only, never fed back into the cell. A probe reports, for a
lightly free-driven pen tip, its height above the touched table and its signed distance to every
registered face, which separates the causes of an early contact. An offline tool fits, to the
recorded contacts of an assembly, the real tube's lateral offset and the real plate's height by grid
search, and reports each mark's distance from the real root.
\figplaceholder{fig:exp_marks}{Photograph of marks on a seam with the measured offsets}{Pen marks on
the weld root and their measured errors.}
\SRC{scripts/touch\_probe.py, contact\_error.py, check\_registration.py, pose\_jitter\_probe.py;
README Gotchas.}

% -------------------------------------------------------------------------------------
\section{Experimental Procedure}
\label{sec:setup_procedure}
One run proceeds as follows.
\begin{enumerate}
\item \textbf{Start.} The robot at its scan home, the External Control program running, the motion
  nodes in dry run until enabled. The calibration chain is checked, not redone: the table height
  under the pen and the camera's view of the table (\S\ref{sec:mc_calibration}).
\item \textbf{Per part.} Place the part; aim the camera by pointing; capture; click the part;
  register; move the part into its place while it is tracked; hold it still for a few seconds; save.
\item \textbf{Refine.} Run the multi-view refinement; record which parts were accepted and the
  camera estimate.
\item \textbf{Compute.} Compute the seams and tacks; run the reachability check.
\item \textbf{Mark.} Plan; execute tack by tack, dots for accuracy; the marking record holds every
  contact.
\item \textbf{Measure.} Measure the marks on the part; archive the run with its assembly.
\end{enumerate}
\TODO{DECIDE the repeated trials (positioning.md): N tacks per assembly (e.g.\ at least 20 per
configuration), single-view and multi-view on the same placement, reported with mean, standard
deviation, median and maximum; and the planned comparisons (FoundationPose vs ICP tracking, APS vs
RRT-Connect, the Cartesian transit vs APS round a tube, bounding boxes vs true shapes, timing per
stage).}

% -------------------------------------------------------------------------------------
\section{WeldSet Benchmark Setup}
\label{sec:setup_weldset}
The benchmark corpus has twelve strata of 60 scenes (\S\ref{sec:ws_corpus}), generated on the
desktop, which holds the authoritative copy of every corpus, the rendered set and the batch results;
the laptop's copies are stale. The batch runs the seven literature methods and the two proposed
extractors over every condition of \S\ref{sec:ws_protocol} in resumable chunks.
\TODO{DESKTOP: corpus sizes and yields per stratum; the final batch (271\,200 rows; about 10.5~h on
the desktop); the CPU; the pinned versions (NumPy, trimesh, Python), since the content hashes depend
on them; the rendered set train\_v1 and the Isaac Sim version.}
\SRC{weld\_generator MANUAL.md \S7, \S8.3; scripts/run\_phase4\_batch.py; memory: project state.}
